\subsection{THETA3-C19: intrinsic CM derivative bounds and the reduced theta frame}\label{C19}

This subsection removes the exact Lubin--Tate coordinate hypothesis
from the derivative conclusions of C18 for a good-reduction CM
elliptic curve at an inert prime.  The argument uses integral
translation operators and the CM action, rather than an asserted
integral isomorphism with a specified polynomial formal group.

\paragraph{Hypotheses and normalization.}
Let $p\geq5$ be inert in an imaginary quadratic field $K$, put
$q=p^2$ and $N=q-1$, and fix the embedding $K\hookrightarrow\C_p$
through which the CM action acts on invariant differentials.
Let $E$ be an elliptic curve over a number field with CM by an order
$\mathcal O\subset K$ whose conductor is prime to $p$.
After a finite extension, suppose that $E$ has good reduction over
$R=\mathcal O_L$, where $L/\Q_p$ is finite and contains $K_p$,
that all its CM endomorphisms are defined, and that its chosen
invariant differential $\omega$ generates the integral invariant
differentials of the smooth model $\mathcal E/R$.
CM endomorphisms extend to this abelian scheme by the N\'eron
mapping property: see Conrad, \emph{Minimal models for elliptic
curves}, author version dated 21 November 2015,
Definition~5.1 and Example~5.3, p.~13.
Uniqueness of the extension preserves the group-homomorphism
identities, since they already hold on the generic fiber.  Thus
$\mathcal O\otimes\Z_p=\mathcal O_{K_p}$, with residue field
$\mathbb F_q$, acts on its formal completion.
The completion of this action is proved below, so no continuity
assumption about an unspecified action is needed.

Choose an integral formal parameter $t$ with
$\omega=(1+O(t))\,dt$, and let $D$ be the invariant derivation dual
to $\omega$.  Thus $Dt(0)=1$.  All valuations below satisfy
$v_p(p)=1$.  The ring $R$ may be ramified over $\Z_p$; in
particular, $R/pR$ is not assumed reduced.

\paragraph{Statements.}
For every $f\in R[[t]]$ and every $n\geq1$,
\begin{equation}\label{C19-strong}
 v_p\bigl(D^nf(0)\bigr)
 \geq\max\left\{0,
       \left\lfloor\frac{np}{p^2-1}\right\rfloor-1\right\}.
\end{equation}
The loss of one cannot be replaced uniformly by zero: at
$n=p^2-1$, the function $f=t^{p^2-1}$ has
$v_p(D^nf(0))=p-1$, whereas
$\lfloor np/(p^2-1)\rfloor=p$.
If the torsion-translation trace of $f$ is zero, the stronger bound
\begin{equation}\label{C19-trace-zero}
 v_p\bigl(D^nf(0)\bigr)
       \geq\left\lfloor\frac{np}{p^2-1}\right\rfloor
\end{equation}
holds.  Independently, the full integral operator satisfies
\begin{equation}\label{C19-Dp}
 D^p(R[[t]])\subset pR[[t]].
\end{equation}

For the theta assertion use a short Weierstrass equation and
differential defined over a number field and integral at the
chosen place.  Write the equation as
$y^2=4x^3-g_2x-g_3$ with unit discriminant,
$\omega=dx/y$, and $t=-2x/y$.  Let $\Lambda$ be the complex period
lattice for this algebraic differential, and let $e_2^*(\Lambda)$
be the coefficient in the classical sigma quasiperiod formula.
Then this coefficient is algebraic and
\begin{equation}\label{C19-e2-integral}
 e_2^*(\Lambda)\in R.
\end{equation}
For the tame algebraic sigma translates in C18, put
\[
 s_T^{\rm red}(t)=
 \exp\bigl(-e_2^*(\Lambda)\lambda(t)^2/2\bigr)s_T(t).
\]
These are the reduced, or Fock-translated, theta normalizations
with the same torsion lifts as C18.  Every finite product and
quotient of these series of total integer tensor degree $d$ has
the form $Q(t)^d v(t)$ with $v\in R[[t]]^\times$ and
\[
 Q(t)=\exp\bigl(-e_2^*(\Lambda)\lambda(t)^2/2\bigr)
            \frac{\sigma(\lambda(t))}{t}.
\]
For every $v\in R[[t]]$, every integer $d$, and $n\geq1$,
\begin{equation}\label{C19-theta-weak}
 v_p\bigl(D^n(Q^dv)(0)\bigr)
 \geq\max\left\{0,
       \frac{n}{p^2-1}-\frac{s_p(n)}{p-1}\right\}.
\end{equation}
Here $s_p(n)$ is the base-$p$ digit sum.  In degree zero the common
factor cancels, and \eqref{C19-strong} applies directly.
These are respectively slopes $1/(p^2-1)$ and $p/(p^2-1)$;
no equality of the two estimates is asserted.

\paragraph{Proof of the operator congruence over $R/pR$.}
Let $S=R/pR$, and reduce the formal group and $D$ to $S$.
A continuous $S$-linear derivation of $S[[t]]$ is of the form
$b(t)\,d/dt$.  Translation invariance identifies such a derivation
with its tangent value: if $D=w(t)\,d/dt$ and $w(0)=1$, the
invariance identity for the group law gives
$b(t)=b(0)w(t)$.  Thus invariant derivations form the free
rank-one $S$-module generated by $D$.
In characteristic $p$, the iterated Leibniz formula shows that
$D^p$ is a derivation, since its middle binomial coefficients
vanish.  It is still translation invariant.  Consequently
$D^p=H D$ for some $H\in S$.

Choose $a\in\mathcal O$ whose image in $\mathbb F_q$ lies outside
$\mathbb F_p$.  The formal endomorphism $[a]$ has unit tangent
coefficient $a$, and so is a formal automorphism.  The invariant
differential identity $[a]^*\omega=a\omega$ gives
$D[a]^*=a[a]^*D$.  Comparing its $p$-fold iterate with
$D^p=HD$ gives $(a^p-a)H=0$ in $S$.
The element $a^p-a$ has nonzero residue and is a unit in $S$,
including when $S$ is nonreduced.  Hence $H=0$, proving
\eqref{C19-Dp}.  Iteration also gives
\begin{equation}\label{C19-iterated-Dp}
 D^k(R[[t]])\subset p^{\lfloor k/p\rfloor}R[[t]]
       \qquad(k\geq0).
\end{equation}

\paragraph{Completion and linearization of the CM action.}
For completeness, let $I=(p,t)\subset R[[t]]$.
For every integral one-dimensional formal group,
$[p](t)=pt+O(t^2)\in I^2$ and
$[p^r](t)\in I^{r+1}$, by induction: substitution of
$u\in I^m$ into $[p](u)$ gives an element of $I^{m+1}$.
If $a-b\in p^r\mathcal O$, then the formal group identity
$[a]=F([b],[a-b])$ shows that
$[a](t)-[b](t)\in I^{r+1}$.
Approximating an element of $\mathcal O_{K_p}$ by elements of
$\mathcal O$ therefore gives a convergent sequence of formal
endomorphisms.  The ring $R[[t]]$ is $I$-adically complete;
for each fixed coefficient this convergence is $p$-adic.
The limit respects the group law, addition, composition, and
tangent coefficients.  This constructs the asserted continuous
$\mathcal O_{K_p}$ action.

Let $\mu_N\subset\mathcal O_{K_p}^\times$ be the Teichm\"uller
roots of unity.  The integral parameter
\[
 X(t)=\frac1N\sum_{\zeta\in\mu_N}\zeta^{-1}[\zeta](t)
\]
has linear coefficient one and satisfies
$X([\eta](t))=\eta X(t)$ for every $\eta\in\mu_N$.
It is thus an invertible integral change of parameter, in which
$[\zeta](X)=\zeta X$.
Let $L(X)$ be the formal logarithm normalized by $\omega$.
Its identity $L(\zeta X)=\zeta L(X)$ implies that only powers
congruent to one modulo $N$ occur.  In particular,
\begin{equation}\label{C19-log-gap}
 L(X)=X+O(X^q).
\end{equation}

Define integral translation operators $D^{[j]}$ by
\[
 f(F(X,Y))=\sum_{j\geq0}(D^{[j]}f)(X)Y^j.
\]
They preserve $R[[X]]$, are $R$-linear and $p$-adically continuous,
and commute with one another.  The last assertion follows by
comparing the coefficients in
$f(F(F(X,Y),Z))$, which is symmetric in $Y,Z$.
Over $L$, the invariant differential Taylor formula is
\[
 f(F(X,Y))=\sum_{k\geq0}\frac{D^kf(X)}{k!}L(Y)^k.
\]
Therefore \eqref{C19-log-gap} gives
\begin{equation}\label{C19-small-Hasse}
 D^{[j]}=D^j/j!\qquad(0\leq j<q).
\end{equation}
No such identification is used for the higher-index operators.

\paragraph{The torsion power sums.}
The generic $p$-torsion is a one-dimensional vector space over
$\mathcal O/p\mathcal O=\mathbb F_q$.  The special fiber's
geometric $p$-torsion is also an $\mathbb F_q$-module, but it has
at most $p$ points: multiplication by $p$ has degree $p^2$ and
inseparable degree at least $p$, since its differential is zero.
It must consequently have only the identity point.
Thus all the generic $p$-torsion reduces to the identity, and
the formal multiplication map has degree $q$.
For any nonzero torsion point with coordinate $\xi$, all its
nonzero torsion coordinates are $\zeta\xi$, $\zeta\in\mu_N$.

Weierstrass preparation, followed by this description of the
roots over characteristic zero, gives
\[
 [p](X)=U(X)(X^q-\beta X),\qquad
 U\in R[[X]]^\times,\quad
 \beta=\xi^N=-p/\alpha,\quad \alpha=U(0)\in R^\times.
\]
Indeed, the prepared polynomial is monic of degree $q$, and its
simple roots are exactly the $q$ torsion coordinates; comparison
of the derivative at zero gives $-U(0)\beta=p$.
Consequently the torsion power sums are exactly
\begin{equation}\label{C19-power-sums}
 \sum_{u\in E[p]}X(u)^j=
 \begin{cases}
 q,&j=0,\\
 0,&j>0\text{ and }N\nmid j,\\
 N(-p/\alpha)^r,&j=rN, r\geq1.
 \end{cases}
\end{equation}
Only the torsion polynomial has this special form; no assertion
that $[p](X)$ itself equals $pX+\alpha X^q$ has been made.

\paragraph{The trace-zero estimate without a polynomial coordinate.}
Define the torsion-translation trace after a finite extension
containing the torsion by
\[
 \mathcal T f(X)=\sum_{u\in E[p]}f(F(X,X(u))).
\]
All $X(u)$ lie in the maximal ideal.  Taylor substitution and
summation thus converge coefficientwise $p$-adically for every
integral $f$.  If $\mathcal T f=0$, equations
\eqref{C19-power-sums} and \eqref{C19-small-Hasse} give
\[
 0=p^2f+N\sum_{r\geq1}(-p/\alpha)^rD^{[rN]}f,
 \qquad
 D^{[N]}f=\frac{\alpha p}{N}(1+H)f,
\]
where
\begin{equation}\label{C19-H}
 H=N\sum_{r\geq2}(-1)^r p^{r-2}\alpha^{-r}D^{[rN]}.
\end{equation}
This series is an integral continuous operator on $R[[X]]$:
each summand preserves the integral ring and its scalar has
valuation tending to infinity.  The ring is complete for this
$p$-adic topology.  Every translation operator commutes with
$D$ and with $\mathcal T$; hence $H$ does also, and all these
operators preserve the trace-zero submodule.

Since $v_p(N!)=p-1$, multiplication by $N!$ gives on that
submodule
\[
 D^N=p^p u(1+H),\qquad
 u=\frac{\alpha N!}{Np^{p-1}}\in R^\times.
\]
Write $n=mN+r$ with $0\leq r<N$.  The commutation and
trace-zero stability just proved allow iteration, giving
\[
 D^nf=p^{pm}u^m r!D^{[r]}(1+H)^m f.
\]
The last factor is integral, so its value at zero has valuation
at least zero.  Furthermore
\[
 v_p(r!)=\left\lfloor\frac rp\right\rfloor
        =\left\lfloor\frac{rp}{p^2-1}\right\rfloor
          \qquad(0\leq r<p^2-1).
\]
For the second equality, write $r=ap+b$, $0\leq b<p$;
then $rp=a(p^2-1)+(bp+a)$, and
$0\leq bp+a<p^2-1$ in the stated range.
This proves \eqref{C19-trace-zero}.

This is the divisibility argument of Katz,
\emph{Divisibilities, congruences, and Cartier duality},
\S III, equations (3.1)--(3.8), pp.~671--672, and Part I of
the proof of Theorem~3.11, pp.~672--673, applied using
\eqref{C19-power-sums} rather than an exact polynomial
coordinate.  The stronger Cartier-operator congruences used
for the other conclusions of that theorem are not claimed here.

\paragraph{The trace ideal and removal of the trace-zero condition.}
Put $A=R[[S]]$, $B=R[[X]]$, with $S=[p](X)$.
This is a finite free extension of degree $q$.
Weierstrass preparation gives
\[
 [p](X)-S=U_S(X)P_S(X),
\]
where $P_S$ is monic of degree $q$ and $U_S$ is a unit.
Write $\omega=a(X)dX$ with $a(X)$ a unit.
The invariant differential identity implies
\[
 [p]'(X)=p\,\frac{a(X)}{a([p](X))}.
\]
In $B=A[X]/(P_S)$ it follows that
$P_S'(X)=p\,v(X)$ with $v\in B^\times$.

Here is an elementary proof of the exact trace ideal, avoiding
an uncomputed different.  If $\ell_{
\rm top}:B\to A$ takes the coefficient of $X^{q-1}$ in the
unique representative of degree less than $q$, then
\[
 \operatorname{Tr}_{B/A}\bigl(b/P_S'(X)\bigr)
      =\ell_{\rm top}(b).
\]
This identity follows over the fraction field by Lagrange
interpolation at the distinct roots of $P_S$; it then holds
in the fraction algebra for every $b\in B$.
Substituting $vb$ gives
\[
 \operatorname{Tr}_{B/A}(b)/p=\ell_{\rm top}(vb)\in A.
\]
Taking $b=v^{-1}X^{q-1}$ makes the right side one.
We have therefore proved the exact equality
\begin{equation}\label{C19-trace-ideal}
 \operatorname{Tr}_{B/A}(B)=pA.
\end{equation}

For any integral $f$, set
$h=\operatorname{Tr}_{B/A}(f)/p\in A$ and
$g=pf-[p]^*h$.  Since the rank is $q=p^2$,
$\operatorname{Tr}(g)=p\operatorname{Tr}(f)-qh=0$.
The algebra trace, pulled back to $B$, equals $\mathcal T$:
after adjoining the torsion, the roots of
$[p](Y)=[p](X)$ are precisely $F(X,X(u))$.
The equality can first be checked in the separable fraction
algebra and then in the integral power-series ring.
Thus \eqref{C19-trace-zero} applies to $g$.
On the other hand, invariant differentiation gives
\[
 D^n[p]^*h(0)=p^nD^nh(0),
\]
where on the right $D$ is the same invariant derivation in the
target parameter.  Since $D$ preserves the integral ring and
$\lfloor np/(p^2-1)\rfloor\leq n$, the identity
$pf=g+[p]^*h$ proves \eqref{C19-strong}.
The additional lower bound zero follows from integrality of
$D^nf$.  Finally, $D=w(t)d/dt$ with $w(0)=1$ gives
$D^n(t^n)(0)=n!$: terms with fewer than $n$ ordinary derivatives
still vanish at zero.  For $n=p^2-1$, Legendre's formula gives
$v_p(n!)=p-1$, proving the claimed sharpness.
The estimates also extend to $f\in\mathcal O_{\C_p}[[t]]$
on the base-changed fixed formal group.  For fixed $n$, the
value $D^nf(0)$ depends only on the first $n+1$ coefficients
of $f$.  Approximate these by integral algebraic coefficients,
apply the finite-extension assertion, and use continuity and
closedness of each stated valuation inequality.

\paragraph{Integral de Rham CM splitting and $e_2^*$.}
Use the short Weierstrass model specified above and set
$\eta=x\omega$.  C18's integral expansion gives
\[
 x=t^{-2}+O(t^2),\qquad
 \omega=(1+O(t^4))dt,\qquad
 \eta=t^{-2}dt+O(t^2)dt.
\]
On $\mathcal E\setminus O$ the differential $\eta$ is regular.
Near $O$ the differential $\eta-d(-1/t)$ is regular.  These
two differentials and the transition function $-1/t$ define
an integral \v Cech--de Rham class.  The image of this class
in $H^1(\mathcal E,\mathcal O_{\mathcal E})$ pairs with
$\omega$ by the residue $\pm\operatorname{Res}_O(\omega/t)
=\pm1$.  Relative Serre duality and the Hodge exact sequence
therefore show that $(\omega,\eta)$ is an integral basis of
$H^1_{\rm dR}(\mathcal E/R)$.  Equivalently, its de Rham cup
pairing is a unit; the sign of this pairing is immaterial here.
For a curve, the Hodge exact sequence used here follows also
directly from the two-term \v Cech--de Rham complex: the map
$H^1(\mathcal O)\to H^1(\Omega^1)$ induced by $d$ is zero,
as the residue of an exact differential is zero.

Choose the global CM element $a$ above.  Both $a$ and
$a-\bar a$ are units in $R$, since reduction of complex
conjugation on $K_p$ is the $p$-power automorphism of
$\mathbb F_q$.  The integral endomorphism acts as
\[
 [a]^*\omega=a\omega,
 \qquad [a]^*\eta=b\eta+c\omega,
 \qquad b,c\in R.
\]
The cup product is multiplied by $\deg[a]=a\bar a$;
comparison in the integral basis gives $ab=a\bar a$ and
hence $b=\bar a$.

Here is the comparison with the specified analytic coefficient.
Under complex uniformization, $x=\wp(z)$ and $\omega=dz$.
Writing $\zeta=\sigma'/\sigma$, one has
$\zeta'=-\wp$.  Thus the periods of $\eta$ are the negatives
of the sigma quasiperiods.  The classical formula is
\[
 \eta_{\rm per}(\gamma)
    =e_2^*(\Lambda)\gamma+C\bar\gamma
       \qquad(\gamma\in\Lambda),
\]
for the nonzero area constant $C$.
Consequently the periods of
$\eta+e_2^*(\Lambda)\omega$ are $-C\bar\gamma$;
this class is a $\bar a$-eigenvector for $[a]^*$.
Comparison with the integral matrix above gives exactly
\[
 c=(\bar a-a)e_2^*(\Lambda).
\]
It proves algebraicity and \eqref{C19-e2-integral}, since
$\bar a-a$ is a unit.  The source for the quasiperiod formula
and its normalization is Bannai--Kobayashi,
\emph{Algebraic theta functions and the $p$-adic interpolation
of Eisenstein--Kronecker numbers}, arXiv:math/0610163v4,
Example~1.9, pp.~10--11.  That example explicitly gives
$\theta_{\rm red}(z)=e^{-e_2^*z^2/2}\sigma(z)$ and
$\theta_{\rm red}'(0)=1$.

\paragraph{Completion of the reduced theta estimate.}
Let $q_0(t)=\sigma(\lambda(t))/t$; this is the common factor
denoted $q(t)$ in C18 and is distinct from the number $q=p^2$.
C18 proves $s_T=q_0u_T$, with $u_T$ an integral unit, and
\[
 B=D^2\log q_0=-x-D^2\log t\in t^2R[[t]].
\]
The sigma quasiperiod identity also gives, by direct substitution
in the translation formula of C18,
\[
 s_T^{\rm red}(t)=Q(t)u_T(t),\qquad
 Q(t)=e^{-e_2^*\lambda(t)^2/2}q_0(t).
\]
In particular every stated product or quotient has the announced
factorization.  No integrality of the raw series $Q$ is assumed.

Put $l_j=D^j\log Q(0)$.  In the additive coordinate
$z=\lambda(t)$, the series $Q$ and its logarithm are even,
and $Q(0)=1$.  Hence $l_0=l_1=0$, all odd $l_j$ vanish,
$l_2=-e_2^*\in R$, and for even $j\geq4$,
\[
 l_j=D^{j-2}B(0),\qquad
 v_p(l_j)\geq\left\lfloor\frac{j-2}{p}\right\rfloor
\]
by \eqref{C19-iterated-Dp}.
The Bell-polynomial proof in C18 now applies verbatim, allowing
the additional even index $j=2$.  Its required numerical
inequality for that index is
\[
 v_p(2!)-0+\frac1{p-1}
       \leq\frac{2p}{p^2-1},
\]
which holds for $p\geq5$.  For clarity, in each exponential
Bell term the inequality
\[
 v_p(j!)-\left\lfloor\frac{j-2}{p}\right\rfloor
       +\frac1{p-1}\leq\frac{jp}{p^2-1}
       \qquad(j\geq2\text{ even})
\]
and $v_p(m!)\leq m/(p-1)$ bound the total denominator loss
by $np/(p^2-1)$.  All $D^n(Q^d)(0)$ are also integral,
since exponential Bell polynomials have integer coefficients
in the integral variables $d l_j$, including for $d<0$.
Thus
\[
 v_p(D^n(Q^d)(0))
       \geq\max\{0,v_p(n!)-np/(p^2-1)\}.
\]
For $v\in R[[t]]$, \eqref{C19-iterated-Dp} bounds the
coefficient of $z^a$ in $v(t(z))$ below by
$\lfloor a/p\rfloor-v_p(a!)\geq-a/(p(p-1))
\geq-ap/(p^2-1)$.
Convolution gives the same factorial bound for $Q^dv$,
and Leibniz's rule gives integrality of its derivatives.
Legendre's formula yields \eqref{C19-theta-weak}.
The degree-zero assertion follows from cancellation of $Q$
and the already proved \eqref{C19-strong}.
\hfill$\square$

\paragraph{Formal-module generalization and several independent factors.}
The formal-group assertions have a slightly wider scope than
the elliptic application.  Let $R$ be any complete discrete
valuation ring of mixed characteristic $(0,p)$ containing
the integers of the unramified quadratic extension of $\Q_p$.
Let $G/R$ be a one-dimensional smooth commutative formal group
with a strict action of those integers, meaning that the action
of $a$ on the tangent space is multiplication by $a$.
Assume that $[p]:G\to G$ is finite flat of degree $p^2$ and
that its reduction has its entire kernel supported at the
identity.  Equivalently in this setting one may assume that
the closed-fiber formal group has height two.
Then all of \eqref{C19-Dp}, \eqref{C19-trace-zero},
\eqref{C19-trace-ideal}, and \eqref{C19-strong} hold.
For these formal assertions alone $p\geq2$ suffices.

Indeed, the integral action is now supplied as a hypothesis,
and averaging over its Teichm\"uller units gives the same
coordinate $X$.  The prepared $p$-torsion polynomial has
$p^2$ simple roots in characteristic zero, all in the maximal
ideal.  Its roots form a one-dimensional vector space over
$\mathbb F_{p^2}$, so the same orbit and power-sum calculation
applies.  Every subsequent translation-operator, convergence,
invariant-differential and trace argument was purely formal.
The elementary factorial identities remain valid at $p=2,3$.
No elliptic sigma or de Rham assertion at those two primes
is included in this generalization.

The proof in fact gives operator divisibility, not only a
bound at the origin.  With
\[
 b(n)=\max\{0,\lfloor np/(p^2-1)\rfloor-1\},
\]
one has $D^nR[[t]]\subset p^{b(n)}R[[t]]$.
For a product $G_1\times\cdots\times G_r$ of such formal
modules and its independently normalized derivations $D_i$,
it follows for every integral series $f$ that
\[
 v_p\bigl(D_1^{n_1}\cdots D_r^{n_r}f(0)\bigr)
       \geq\sum_{i=1}^r b(n_i).
\]
To justify coefficient rings in this last assertion, apply the
one-variable proof coefficientwise in the other variables:
all translation operators preserve their integral coefficients,
the same $p$-adically convergent sums still converge, and the
finite-free trace identity is preserved by this power-series
base change.  The $D_i$ commute, so their integral operator
divisibilities multiply.

Thus these bounds apply whenever the formal completion of a
good-reduction CM abelian scheme has been decomposed into the
specified one-dimensional height-two factors by its integral
local CM idempotents.  The precise decomposition, its integral
differentials and the theta-section connection are still part
of the input when using this multi-factor consequence.

\paragraph{Scope of the arithmetic bridge.}
The only formal-group hypothesis used here beyond good reduction
is the maximal local unramified quadratic CM action; a CM order
with conductor divisible by $p$ is not covered by this proof.
The arguments apply after every finite extension carrying the
stated good model and normalized differential.  In particular,
they never replace $R/pR$ by its reduced quotient.
The result supplies intrinsic derivative estimates for the actual
tame theta functions with the named differential and torsion
normalization.  A comparison with a differently normalized
harmonic or covariant derivative still uses its explicit analytic
normalization; the bounds do not make that identification by
the existence of an arbitrary line-bundle trivialization.

The de Rham part is a concrete realization of the integral
correspondence and polarization setting, and its algebraic
splitting is computed rather than assumed.  It is compatible
with the Albanese and polarization discussion in Grothendieck's
cote~44, scans~7--8, and with C17's precisely stated cycle-class
specialization formalism.  Those archive passages are not cited
as proofs of the local valuation estimates established here.
